\chapter{Experimentos numéricos y comparación de modelos}\label{ch:numerics}
\objetivos{Reproducir cuatro familias de práctica con ejemplos nuevos; analizar convergencia y discrepancias; diseñar una adquisición de datos sin inventar ensayos.}

\section{Organización de los experimentos}
Las cuatro familias siguen las competencias de las prácticas facilitadas: vuelo horizontal, ascenso, consumo en crucero y comparación de configuraciones. Se utilizan condiciones originales y se proporciona una metodología. No son una memoria de prácticas del estudiante ni se atribuyen ensayos personales o datos de SAULE inexistentes.

Cada experimento conserva entradas, salidas con unidades, registro de ejecución y limitaciones. Los gráficos del libro se reconstruyen a partir de los CSV o de las ecuaciones indicadas. Un gráfico sintético lleva esa identificación; uno procedente de JSBSim indica el modelo y condiciones.

\section{Familia 1: barrido de vuelo horizontal}
La rejilla original del T-6 combina masas nominales de 5.600 y 5.850 lbm, altitudes de 5.000 y 9.000 ft y once CAS entre 115 y 215 kt. Se obtienen 44 trimas longitudinales. Se registran EAS y TAS resultantes: el eje de resistencia se dibuja frente a EAS, aunque la condición de entrada sea CAS.

En una polar ideal, a igual masa y configuración, la resistencia frente a EAS depende poco de altitud si los coeficientes permanecen constantes. En un modelo completo pueden aparecer variaciones por Mach, Reynolds, trim, empuje inclinado y dispositivos. La coincidencia aproximada no debe imponerse por construcción.

\ilustra{jsbsim-t6}{Rejilla de trim realmente ejecutada con el modelo educativo t6texan2. Las curvas se trazan frente a EAS resultante; cada punto procede de una instancia inicializada de nuevo.}

La potencia necesaria de empuje se obtiene como $DV_{\rm TAS}$. A igual EAS y masa, la TAS aumenta con altura; por tanto, aun con resistencia aproximadamente igual, la potencia de empuje necesaria aumenta. Dibujar $DV_{\rm EAS}$ ocultaría precisamente ese efecto.

\section{Familia 2: estimación de ascenso}
Para cada punto de la rejilla se estabiliza el motor a mando máximo sin avanzar la trayectoria y se consulta el empuje disponible. Se calcula
\begin{equation}
P_d=T_{\max}V_{\rm TAS},\quad P_n=DV_{\rm TAS},\quad
\dot h_{\rm aprox}=\frac{P_d-P_n}{W},\quad
\sin\gamma_{\rm aprox}=\frac{T_{\max}-D}{W}.
\end{equation}
Es una estimación de actuaciones a partir de trimas horizontales, con $L\simeq W$; no se presenta como trim exacto de ascenso. La máxima tasa y el máximo ángulo se localizan con objetivos distintos. Con once puntos de velocidad sólo se tiene un máximo de rejilla; para afinar Vx o Vy debe refinarse el intervalo y comprobar convergencia en velocidad.

\ilustra{jsbsim-ascenso}{Estimación de ascenso por exceso de potencia para una de las condiciones originales del T-6. Se utiliza potencia de empuje, de acuerdo con el tipo de motor y propulsor realmente configurados.}

La tabla de motor puede usar Mach y altitud de densidad en pies. Leerla como altitud geométrica en metros produce una consulta incorrecta aunque la interpolación funcione. Multiplicar un empuje por un coeficiente tabulado sigue dando empuje: la velocidad TAS es necesaria para convertirlo en potencia.

\section{Familia 3: crucero del 747 con consumo}
Se inicializa el modelo B747 a 27.000 ft y CAS de 260 kt, con 36.000 lbm de combustible repartido entre los cuatro depósitos alimentadores y el quinto vacío. Se simulan 1.800 s con paso de 0.05 s y un control docente explícito de altura y TAS. No se impone masa constante.

El control usa error de velocidad y su integral para gases y error de altura, velocidad vertical, cabeceo y tasa de cabeceo para elevador. Añade amortiguación lateral sencilla. Los mandos se limitan al intervalo permitido. Se registran errores máximos de seguimiento y no se describe el vuelo como altitud y velocidad exactamente constantes si los datos muestran deriva.

\ilustra{jsbsim-b747}{Combustible calculado durante el crucero original de 30 minutos. El consumo procede de sumar los cuatro motores; no se presenta como dato experimental de un Boeing real.}

\subsection{Comparaciones analíticas honestas}
Una primera referencia congela el gasto inicial y calcula $\Delta m=\dot m_{f,0}t$. Es una aproximación de gasto constante, distinta de la integral de masa variable. Otra comparación identifica una polar y un consumo efectivo locales y usa la integral a TAS y densidad constantes. Sólo es legítimo llamar a esa segunda comparación modelo equivalente si se documenta el ajuste y su intervalo; no se atribuyen sus parámetros a un manual del fabricante.

La masa de combustible se contrasta con $\int\sum_i\dot m_{f,i}\,\dd t$. El peso total se convierte a masa y se compara con masa vacía más combustible. Puede existir un desfase pequeño de actualización de una propiedad respecto a otra dentro del paso; se mide y registra en vez de interpretarlo inmediatamente como creación de masa.

La ventana de media hora se complementa con dos simulaciones completas hasta agotar los depósitos alimentadores: 28\,000 y 36\,000 lbm iniciales, con el quinto depósito vacío. Se conserva el mismo esquema docente de control. La terminación se identifica mediante el contenido total accesible; el último paso tiene resolución temporal de 0.1 s.

\begin{center}\begin{tabular}{rrrr}\toprule
Combustible inicial (lbm)&Tiempo simulado (s)&Modelo equivalente (s)&Diferencia\\\midrule
28000&4113.0&4105.3&0.187\\
36000&5264.4&5313.6&-0.926\\
\bottomrule\end{tabular}\end{center}
El modelo equivalente ajusta $D=A_V+B_VW^2$ a cinco trimas del mismo modelo de avión y usa un consumo específico efectivo inicial. No es una fuente independiente ni una identificación experimental. El ajuste de resistencia tiene un error relativo máximo inferior al 0.01\% en esos cinco puntos; no describe por ello cualquier condición de vuelo.

Las alturas derivan hasta aproximadamente 58 y 75 m respecto al objetivo en las dos autonomías, aunque el error máximo de TAS permanece por debajo de 0.06 m/s. Por tanto, la comparación no cumple exactamente altitud constante. El consumo inicial registrado también puede diferir ligeramente del que se establece al avanzar el motor. Esas limitaciones y la ausencia de un estudio de refinamiento de paso para las autonomías completas se conservan en el informe; la pequeña diferencia numérica no se presenta como validación experimental.

\section{Convergencia temporal en un caso libre}
El T-6 parte del mismo trim a CAS 155 kt y 5.000 ft y se integra libremente durante 60 s con pasos 0.10, 0.05 y 0.025 s. Se comparan altura, TAS, masa, actitud y aceleraciones finales. La diferencia entre resultados sucesivos aporta una estimación empírica de sensibilidad al paso.

No se obtiene un orden formal sólo porque dos diferencias decrezcan. Los estados de motor, actualización de combustible, tolerancia de trim y términos terrestres pueden limitar el régimen asintótico. Se exige convergencia de las magnitudes de interés para ese caso, y se conserva como alcance separado de la convergencia de la ODE analítica de crucero.

\section{Familia 4: método para SAULE sin datos inventados}
No se ha accedido al simulador físico SAULE. Se propone un protocolo: fijar masa, combustible, potencia, hélice, temperatura y configuración; establecer un intervalo de altitud común; estabilizar velocidad indicada; registrar tiempo, altitud, IAS y datos de viento y temperatura disponibles; repetir con flaps retraídos y con una configuración de flaps definida. No se atribuye una IAS a TAS sin las correcciones necesarias.

La tasa media se calcula por ajuste de $h(t)$ en el intervalo estabilizado, no diferenciando directamente una señal ruidosa. El ángulo requiere velocidad horizontal aérea o terrestre según lo que se quiera comparar:
\begin{equation}
\sin\gamma_a=\frac{\dot h-U_h}{V_{\rm TAS}},\qquad
\tan\gamma_g=\frac{\dot h}{V_{g,H}}.
\end{equation}
Para un tramo de altura limitada, se reporta dispersión entre repeticiones y se descartan transitorios mediante un criterio previo. Elegir después sólo el intervalo que más favorece una configuración sesga la comparación.

\subsection{Qué se puede predecir sobre flaps}
Los flaps suelen modificar $C_{L\max}$, polar, trim y velocidad óptima. A una misma TAS y peso, la comparación de ángulo depende de empuje disponible y resistencia de cada configuración. Un aumento de sustentación máxima permite menor velocidad, pero no garantiza mayor ángulo de ascenso: el incremento de resistencia puede dominar. Para decidir se compara el máximo de $T_d-D$ dentro del dominio admisible de cada configuración, con su rendimiento de hélice y límites.

Si se incluyen datos sintéticos para practicar un ajuste, deben etiquetarse en cabecera, gráfico y texto. No se mezclan en una tabla con mediciones sin una columna de procedencia. Esta edición entrega el método y no presenta ensayos SAULE como ejecutados.

\section{Incertidumbre y discrepancia}
Una diferencia entre teoría y simulación puede proceder de error numérico, definición de magnitud, aproximaciones físicas o parámetros. Se revisan primero unidades, ejes y condiciones, después discretización y finalmente modelo. Ajustar una polar no corrige una conversión de nudos ni una suma incompleta de motores.

La incertidumbre de parámetros puede propagarse por sensibilidad: se perturba un parámetro dentro de su intervalo y se recalcula la magnitud. Esa variación no equivale a una distribución estadística sin justificar probabilidades. Para modelos educativos con parámetros estimados, un intervalo de sensibilidad suele ser más honesto que una barra de error aparentemente experimental.

\section{Taller resuelto: Interpretar un refinamiento realmente ejecutado}

\noindent\textbf{Nivel:} Intermedio.

\apartado{1. Enunciado}{En el ensayo libre de 60 s del T-6 se usaron pasos 0.10, 0.05 y 0.025 s. Compare las TAS finales y calcule diferencias sucesivas. ¿Permiten asignar un orden universal a JSBSim?}

\apartado{2. Datos}{Los tres CSV y registros de versión están incluidos en el proyecto; las condiciones iniciales y duración son comunes.}

\apartado{3. Incógnitas}{Las magnitudes y comprobaciones solicitadas en el caso.}

\apartado{4. Sistema de referencia}{Misma propiedad TAS SI al final de cada trayectoria.}

\apartado{5. Hipótesis físicas}{Comparación de un experimento libre concreto, sin atribuir todo el error a una sola potencia de paso.}

\apartado{6. Modelo matemático}{Interpretar un refinamiento realmente ejecutado}

\apartado{7. Ecuaciones simbólicas}{Se parte de los balances y definiciones siguientes.}

\begin{equation}\Delta_1=|V_{0.10}-V_{0.05}|\end{equation}

\begin{equation}\Delta_2=|V_{0.05}-V_{0.025}|\end{equation}

\begin{equation}p_{observado}=\ln(\Delta_1/\Delta_2)/\ln2\end{equation}

\apartado{8. Desarrollo algebraico}{Las operaciones se realizan en este orden, conservando las condiciones del modelo.}

\begin{enumerate}

\item Se extraen los valores finales sin redondear antes de restarlos.

\item Se calculan diferencias y un orden observado sólo como indicador condicionado.

\item La mezcla de actualización de masa, tolerancia de trim e integradores impide generalizarlo a todos los estados y casos.

\end{enumerate}

\apartado{9. Sustitución numérica}{Se usan las magnitudes del apartado de datos en SI; los ángulos se convierten a radianes antes de evaluar funciones trigonométricas. Las densidades indicadas como ISA se obtienen de la capa correspondiente.}

\noindent Diferencia primera:\quad $5.151754832866118 \cdot 10^{-5}\simeq 5.151755\times10^{-5}\,\mathrm{m\,s^{-1}}$.\par

\noindent Diferencia segunda:\quad $2.087958689855895 \cdot 10^{-5}\simeq 2.087959\times10^{-5}\,\mathrm{m\,s^{-1}}$.\par

\noindent Orden observado condicionado:\quad $1.302970770044629\simeq 1.302971\,$.\par

\apartado{10. Resultados con unidades}{Los resultados pertenecen al modelo y al intervalo de datos declarados.}

\begin{center}\begin{tabular}{lr}\toprule Magnitud&Resultado\\\midrule

Diferencia primera & $5.151755\times10^{-5}\,\mathrm{m\,s^{-1}}$\\

Diferencia segunda & $2.087959\times10^{-5}\,\mathrm{m\,s^{-1}}$\\

Orden observado condicionado & $1.302971\,$\\

\bottomrule\end{tabular}\end{center}

\apartado{11. Interpretación}{Las diferencias pequeñas y decrecientes aportan sensibilidad temporal para esta TAS final; no prueban un orden formal global.}

\apartado{12. Comprobación independiente y alcance}{Comprobar también altura, actitud y masa final en los registros y repetir con otro paso si la aplicación exige tolerancias más estrictas.}

\apartado{13. Limitaciones}{Las autonomías completas del 747 no se declaran convergentes por este ensayo corto de otra aeronave.}

\par\smallskip\noindent{\footnotesize\textit{Datos y solución originales. Las sustituciones registradas se reproducen con SymPy y Studium; eso no establece validación experimental ni revisión experta.}}


\repaso{¿Una trayectoria convergente de JSBSim prueba que el consumo de un avión real está bien modelado?}{No. La convergencia examina discretización de ese modelo y experimento. La fidelidad del consumo necesita datos independientes y condiciones comparables.}
\resumen{Los experimentos reproducibles conservan sus entradas y resultados y declaran qué se ejecutó. La rejilla de trim, la estimación de ascenso, la trayectoria de crucero y la metodología de ensayo son evidencias diferentes. Las discrepancias se investigan sin convertir simulación en validación experimental.}
\fuentes{Enunciado Ágora 2026/2027, cuatro familias prácticas, consultado para competencias. JSBSim 1.3.3 y resultados originales del proyecto. Los modelos y las condiciones de esta edición se distinguen de las prácticas evaluables y de datos experimentales no disponibles.}
